\clearpage
\thispagestyle{empty}\SourceLeaf{283}
\pdfbookmark[0]{Notices of the Author's Own Papers}{anzeigen-title}
\begin{center}
\vspace*{2in}
{\Large NOTICES\par}
\vspace{0.5in}
OF THE AUTHOR'S OWN\par
\vspace{0.5in}
{\Large PAPERS.\par}
\end{center}
\clearpage
\SourceLeaf{285}
\fancyhead[L]{\small Notices of the Author's Own Papers}
\pdfbookmark[1]{Göttingische gelehrte Anzeigen. April 5, 1813.}{anzeige-1813-april5}
\begin{center}
\rule{0.85\textwidth}{0.3pt}\par
Göttingische gelehrte Anzeigen. April 5, 1813.\par
\rule{0.85\textwidth}{0.3pt}
\end{center}
\medskip
On March 18, Professor \textsc{Gauss} submitted a lecture to the Royal Society under the title
\begin{center}
\textit{Theoria attractionis corporum sphaeroidicorum ellipticorum homogeneorum, methodo nova tractata.}
\end{center}
As is well known, since \textsc{Newton}'s time the foremost geometers have worked on this famous problem almost as if competing with one another. \textsc{Newton} took the first step by showing how to determine the attraction on a point on the axis of a spheroid of revolution and how to find the simple ratio between the attractions on all points lying on one and the same diameter inside the spheroid. Next, in his famous prize essay on the tides, \textsc{Maclaurin} succeeded in determining the attraction on all points on the surface of the spheroid. By \textsc{Newton}'s theorem the attractions at interior points could also be reduced to these, so that only the attraction on exterior points remained to be determined. This, of course, was the most difficult part of the problem. \textsc{Maclaurin} had already taken a step here as well. He determined the attraction on points in the extensions of the equator and of the axis. \textsc{Maclaurin}'s discoveries were universally admired as masterpieces of synthesis and were for a time regarded as proof that there are cases in which the synthetic method has a decided advantage over
\SourceLeaf{286}the analytical method. \textsc{Lagrange} restored the analytical method to its rightful place by bringing within its reach a problem that had seemed accessible only to synthesis, and by showing with his characteristic dexterity how to obtain all of \textsc{Maclaurin}'s discoveries by analysis. Although this made no new advance on the problem itself, it had to be regarded as a highly important preparation for later work. It was \textsc{Legendre} who succeeded in completing the theory of the attraction of spheroids of revolution. He discovered and proved the beautiful theorem that the attraction of a spheroid on an exterior point has the same direction as the attraction on the same point by a second spheroid whose surface passes through that point, provided that the two generating ellipses have the same focus, and that the former attraction is to the latter as the mass of the first spheroid is to that of the second.

All this concerns spheroids generated by revolving half an ellipse about one axis or the other. But there still remained the much harder problem of determining the attraction of an ellipsoid whose equator is also elliptical, that is, of a body whose every section by a plane gives an ellipse. \textsc{Maclaurin} had already indicated how to determine the attraction on points in the directions of the three principal axes, and \textsc{d'Alembert} and \textsc{Lagrange} had given analytical proofs. Further, by induction \textsc{Legendre} had surmised that his beautiful theorem cited above held generally, though he had been unable to find a rigorous proof. It fell to \textsc{Laplace} to fill this gap and complete the solution of the problem in its full generality (1782). One might suppose that this concluded the investigation. Yet the very fact that several geometers have since taken up the same subject afresh shows that much remained to be desired. The first and most essential thing about a problem is always to solve it at all. But several routes often lead to the same goal. Quite often the goal is first reached by a long, thorny detour. The shortest route, the true and genuine one, is discovered much later. \textsc{Laplace}'s solution is a fine example of the most refined analytical art, but the route by which he arrived at it. is long and arduous, and surely few geometers and
\SourceLeaf{287}astronomers have followed him along it. The use of infinite series whose convergence has not been proved also detracts somewhat from the clarity and conclusiveness of the proof. \textsc{Legendre} did give another solution in 1788, though almost everything we have said against \textsc{Laplace}'s applies to it as well. A competent judge, \textsc{Lagrange}, gives the following verdict on the solutions of those two great analysts (in the \textit{Nouv. Mém. de l'Acad. de Berlin} 1793): ‘\textit{On ne peut regarder leurs solutions que comme des chefs-d'oeuvre d'analyse, mais on peut désirer encore une solution plus directe et plus simple: et les progrès continuels de l'analyse donnent lieu de l'espérer.}’ Since then \textsc{Biot} and \textsc{Plana} have also sought to improve and simplify those two proofs. Valuable as these works are, the solutions must still be counted among the most intricate and subtle applications of analysis.

The author of the present paper had long been convinced that the \textit{genuine} method of solving this famous problem still had to be found. Six months ago he was prompted to examine it more closely. Taking a route entirely different from those taken before, he had the pleasure of arriving at a solution so surprisingly short and simple that its essentials could be set out in two pages. He has of course given a somewhat longer account here. For one thing, he wished to make it intelligible to less practiced readers, to whom these investigations, so interesting for the figure of the Earth, had until now been wholly inaccessible. He already has several proofs that the new solution is fully suited to this purpose. For another, it seemed worthwhile to develop the grounds on which it rests somewhat more fully than the immediate purpose required, since they will often be useful on other occasions too.

We now wish to present the main points of the entire solution as briefly as possible, yet in enough detail for those familiar with the subject. Here we must forgo making ourselves fully intelligible to readers still unfamiliar with investigations of this kind. We must refer them to the full paper itself, which has already been printed and will shortly appear in the Society's second volume of \textit{Commentationes recentiores}.

\SourceLeaf{288}The author begins by establishing six different general theorems by which triple integrals taken over a solid region can be reduced to double integrals taken only over the surface of the body. We give only three of these theorems here, since the others are unnecessary for the present investigation.

Let $\mathrm{d}s$ be an element of the surface of a body of any shape. Let $PQ,PM,PX,PY,PZ$ be straight lines drawn from a point $P$ of this element, respectively normal to the surface and outward, toward the attracted point $M$, and parallel to the three coordinate axes. Further, let $r$ be the distance of the point $M$ from $P$, $MQ$ the angle between $PM$ and $PQ$, $MX$ the angle between $PM$ and $PX$, and $QX$ the angle between $PQ$ and $PX$. Finally, let $\pi$ denote the ratio of the circumference of a circle to its diameter, and $X$ the attraction that the whole body exerts on the point $M$ parallel to the coordinate axis $x$. This gives
\[
\text{I.}\qquad \int\frac{\mathrm{d}s.\cos MQ}{rr}=0\text{ or }=-4\pi
\]
according as $M$ lies outside or inside the body, and
\[
\begin{aligned}
\text{II.}\qquad &\int\frac{\mathrm{d}s.\cos QX}{r}=X\\
\text{III.}\qquad &\int\frac{\mathrm{d}s.\cos MQ.\cos MX}{r}=-X
\end{aligned}
\]
where the integrals are to be taken over the entire surface of the body. We omit the proofs of these theorems here and merely observe that the first two are based on dividing the body into conical elements with their apex at $M$, whereas the third is based on dividing it into prismatic elements parallel to the coordinate axis $x$.

For the surface of an ellipsoid whose three semiaxes are $A,B,C$, the coordinates $x,y,z$ satisfy the equation
\[
\frac{xx}{AA}+\frac{yy}{BB}+\frac{zz}{CC}=1
\]
Further, $\cos QX=\frac{x}{AA\rho}$ if one sets, for brevity,
\[
\BareRoot{(\frac{xx}{A^4}+\frac{yy}{B^4}+\frac{zz}{C^4})}=\rho
\]
If $a,b,c$ denote the coordinates of the point $M$, the result is
\SourceLeaf{289}
\[
\begin{aligned}
r&=\BareRoot{[(a-x)^2+(b-y)^2+(c-z)^2]}\\
\cos MX&=\frac{a-x}{r}\\
\cos MQ&=\frac{1}{\rho r}(\frac{(a-x)x}{AA}+\frac{(b-y)y}{BB}+\frac{(c-z)z}{CC})
\end{aligned}
\]
Two new variables $p,q$ are now introduced, on which $x,y,z$ depend as follows.
\[
\begin{aligned}
x&=A\cos p\\
y&=B\sin p\cos q\\
z&=C\sin p\sin q
\end{aligned}
\]
Thus, to cover the entire surface of the ellipsoid, one must let $p$ range from $0$ to $180^0$ and $q$ from $0$ to $360^0$. Finally, set $X=ABC\xi$ as well. For familiar reasons, $\mathrm{d}s=\mathrm{d}p.\mathrm{d}q.ABC\rho\sin p$. If, for brevity, one sets
\[
\frac{(a-x)x}{AA}+\frac{(b-y)y}{BB}+\frac{(c-z)z}{CC}=\psi
\]
the three theorems above then take the following form.
\[
\begin{aligned}
[1]\qquad &\int\int\frac{\mathrm{d}p.\mathrm{d}q.\sin p.\psi}{r^3}=0\text{ or }=-\frac{4\pi}{ABC}\\
[2]\qquad &\int\int\frac{\mathrm{d}p.\mathrm{d}q.\cos p.\sin p}{r}=A\xi\\
[3]\qquad &\int\int\frac{\mathrm{d}p.\mathrm{d}q.\sin p.\psi(a-x)}{r^3}=-\xi
\end{aligned}
\]
Now regard $A,B,C$ as specified particular values of three variables $\alpha,\beta,\gamma$, which are related so that $\alpha\alpha-\beta\beta$ and $\alpha\alpha-\gamma\gamma$ remain constant. Formula [1] readily leads to the conclusion that as $\alpha$ increases without limit, $\xi$ decreases indefinitely. Differentiating [2] with respect to the variables $\alpha,\beta,\gamma$ and using the variation sign $\delta$, one obtains
\[
\begin{aligned}
\alpha\delta\xi+\xi\delta\alpha&=-\int\int\frac{\mathrm{d}p.\mathrm{d}q.\cos p.\sin p.\delta r}{rr}\\
&=\delta\alpha.\int\int\frac{\mathrm{d}p.\mathrm{d}q.\sin p.x\psi}{r^3}
\end{aligned}
\]
\SourceLeaf{290}or, if one substitutes the value from [3] for $\xi$ here,
\[
\alpha\delta\xi=\delta\alpha.\int\int\frac{\mathrm{d}p.\mathrm{d}q.\sin p.a\psi}{r^3}
\]

\Needspace{11\baselineskip}
Comparing this with [1] gives
\[
[4]\qquad\delta\xi=0
\]
if the point $M$ lies outside the ellipsoid, and
\[
[5]\qquad\delta\xi=-\frac{4\pi a\delta\alpha}{\alpha\alpha\beta\gamma}
\]
if $M$ lies inside.

It follows from [4] that $\xi$ is constant, or that the attraction is proportional to the mass, for all spheroids whose principal sections are ellipses with the same foci, provided that $M$ does not lie inside. Determining the attraction of a spheroid on an exterior point thus reduces to determining the attraction of another spheroid described from the same foci and passing through the attracted point. To determine this attraction, consider the other case, where the attracted point lies inside. Substituting
\[
\begin{aligned}
\beta\beta&=\alpha\alpha-AA+BB\\
\gamma\gamma&=\alpha\alpha-AA+CC
\end{aligned}
\]
in equation [5], also setting $\frac{A}{\alpha}=t$ and replacing the sign $\delta$ with the usual $d$, gives
\[
\xi=\frac{4\pi a}{A^3}\int\frac{tt\mathrm{d}t}{\BareRoot{[(1-(1-\frac{BB}{AA})tt)(1-(1-\frac{CC}{AA})tt)]}}
\]
where the integral must be fixed so that it vanishes for $t=0$ and must then be extended to $t=1$ for the \textit{specified} spheroid. Thus in the same sense one obtains
\[
[6]\qquad X=\frac{4\pi aBC}{AA}\int\frac{tt\mathrm{d}t}{\BareRoot{[(1-(1-\frac{BB}{AA})tt)(1-(1-\frac{CC}{AA})tt)]}}
\]
This formula gives the attraction on all points that do not lie outside. Since it must hold up to the surface itself, and since the attraction on exterior \SourceLeaf{291}points has already been reduced to the attraction on points on the surface, the problem is thereby completely solved. (It scarcely needs pointing out that the attractions parallel to the other two principal axes are obtained simply by interchanging $A,a$ with $B,b$ or with $C,c$.)

Equation [6] further shows that, for an interior point, the attraction of all spheroids that are similar to one another and similarly situated is identical. Thus if one imagines such a spheroid divided into layers bounded by similar ellipsoidal surfaces, it is clear that all layers lying outside the attracted point contribute nothing at all to the attraction, and only the attraction of the spheroidal core whose surface passes through the attracted point remains.

In closing, the author mentions the latest work on the same subject, by Mr. \textsc{Ivory} in the \textit{Philos. Transact.} of 1809. Count \textsc{Laplace} drew it to his attention, and the author only became acquainted with it when his own paper was already completely finished. By a very happy idea, Mr. \textsc{Ivory} reduced the attraction on an exterior point to the attraction on an interior point. His method of determining the attraction on interior points themselves does indeed show great ingenuity and skill. Yet, like \textsc{Laplace}'s solution for exterior points, it is based in part on considering infinite series that do not converge everywhere, and it falls far short of the simplicity one might have wished for. So \textsc{Ivory}'s solution to the problem, considered as a whole, is in fact scarcely less elaborate and intricate than the solutions of \textsc{Laplace} and \textsc{Legendre}. In other respects, that solution and the one reported here rest on quite different grounds and have nothing in common except the use of the two variables denoted above by $p,q$.

\bigskip
\begin{center}
[\textit{Manuscript Note.}]
\end{center}
\medskip
By a method similar to the one presented here, one can also determine the function $V$, that is, the sum of all the particles of a body, each divided by its distance from the point being attracted.
